\chapter{Angular avanzato}

Sappiamo che cosa sono i componenti, come si usa il Router integrato e come
gestire \code{@Input} e \code{@Output}. In questo capitolo vediamo altre
funzionalità importanti del framework: i \textbf{form}, i \textbf{servizi} e
l'iniezione delle dipendenze, la \textbf{protezione delle rotte}, il
\textbf{client HTTP} e i \textbf{signal}.

\section{I form}

I form sono in molti casi il modo preferito per raccogliere l'input
dell'utente. In Angular si possono gestire in due modi:

\begin{itemize}
  \item \term{form guidati dal template}: le associazioni fra i dati del
        componente e il form sono definite \textbf{implicitamente}, tramite
        direttive nel template;
  \item \term{form reattivi}: le associazioni sono definite
        \textbf{esplicitamente} nel componente.
\end{itemize}

\subsection{Form guidati dal template}

\begin{lstlisting}[style=ts, name=login.component.ts]
import { Component } from '@angular/core';
import { FormsModule } from '@angular/forms';

@Component({
  selector: 'app-login',
  standalone: true,
  imports: [FormsModule],
  templateUrl: './login.component.html',
  styleUrl: './login.component.scss'
})
export class LoginComponent {
  user: string = "";
  pass: string = "";
}
\end{lstlisting}

\begin{affianco}[0.52][c]
\codicepiccolo
\begin{lstlisting}[style=html, name=login.component.html]
<h1>Login</h1>
<form>
  <label for="usr">Username:</label>
  <input id="usr" name="usr" [(ngModel)]="user"/>

  <label for="pwd">Password:</label>
  <input id="pwd" name="pwd" type="password" [(ngModel)]="pass"/>
</form>
<p>
  Username is <strong>{{user}}</strong>;
  Password is <strong>{{pass}}</strong>
</p>
\end{lstlisting}
\risultato
\begin{browser}[localhost:4200/login]
  {\Lora\bfseries\LARGE Login\par}\vspace{6pt}
  Username: \wtcampo[1.9cm]{alice}\enspace
  Password: \wtcampo[1.9cm]{$\bullet\bullet\bullet\bullet\bullet$}\par\vspace{6pt}
  Username is \textbf{alice}; Password is \textbf{12345}
\end{browser}
\begin{immagine}[\code{[(ngModel)]}: il campo e la proprietà restano\\
                 sempre allineati, nei due versi.]
\begin{tikzpicture}[
    a/.style={-{Stealth[length=4pt,width=3.5pt]}, line width=0.75pt,
      draw=neutral!45!black}]
  \node[wtnodeplain, font=\FiraCodeNL\scriptsize] (f) at (0,0) {<input id="usr">};
  \node[mvar, minimum width=1.9cm] (u) at (5,0) {"alice"};
  \node[mnome, anchor=south] at (u.north) {user};
  \draw[a] ([yshift=3pt]f.east) -- node[mcod, above]{(ngModelChange)}
    ([yshift=3pt]u.west);
  \draw[a] ([yshift=-3pt]u.west) -- node[mcod, below]{[ngModel]}
    ([yshift=-3pt]f.east);
\end{tikzpicture}
\end{immagine}
\end{affianco}

La direttiva \code{[(ngModel)]="pass"} specifica che il valore di un campo del
form è legato a una proprietà del componente.

\info{La \guillemotleft banana nella scatola\guillemotright}{
  La notazione \code{[()]} è nota scherzosamente come \textit{banana in a box} e
  rappresenta un \term{legame bidirezionale}: combina il property binding
  (parentesi quadre, dati che scendono) con l'event binding (parentesi tonde,
  eventi che salgono).

  Il risultato è che, digitando nel campo, i dati del componente si aggiornano
  automaticamente e il template si aggiorna di conseguenza; e se una proprietà
  cambia nel componente, il campo del form si aggiorna da sé. Per verificarlo
  basta inizializzare \code{user} a un valore diverso dalla stringa vuota.}

\subsection{Form reattivi}

\begin{lstlisting}[style=ts, name=login.component.ts]
import { Component } from '@angular/core';
import { FormControl, FormGroup, ReactiveFormsModule } from '@angular/forms';

@Component({
  selector: 'app-login',
  standalone: true,
  imports: [ReactiveFormsModule],
  templateUrl: './login.component.html',
  styleUrl: './login.component.scss'
})
export class LoginComponent {
  loginForm = new FormGroup({
    user: new FormControl(''),    // l'argomento e' il valore iniziale
    pass: new FormControl('')
  });

  onSubmit() {
    alert(`Username: ${this.loginForm.value.user}
           Password: ${this.loginForm.value.pass}`);
  }
}
\end{lstlisting}

\begin{affianco}[0.48][c]
\codicepiccolo
\begin{lstlisting}[style=html, name=login.component.html]
<h1>Login</h1>
<form [formGroup]="loginForm" (ngSubmit)="onSubmit()">
  <label for="usr">Username:</label>
  <input id="usr" formControlName="user"/>

  <label for="pwd">Password:</label>
  <input id="pwd" formControlName="pass" type="password"/>

  <button type="submit">Submit</button>
</form>
<p>
  Username is <strong>{{loginForm.value.user}}</strong>;
  Password is <strong>{{loginForm.value.pass}}</strong>
</p>
\end{lstlisting}
\risultato
\begin{immagine}
\begin{tikzpicture}[
    tp/.style={mdom},
    cp/.style={mfun, rounded corners=3pt},
    l/.style={mrif, dashed, <->}]
  \node[mnota, font=\Lato\scriptsize\bfseries, anchor=east] at (0,0.65) {template};
  \node[mnota, font=\Lato\scriptsize\bfseries, anchor=west] at (2.9,0.65) {componente};
  \node[tp, anchor=east] (f) at (0,0) {<form>};
  \node[tp, anchor=east] (iu) at (0,-1) {<input id="usr">};
  \node[tp, anchor=east] (ip) at (0,-2) {<input id="pwd">};
  \node[cp, anchor=west] (g) at (2.9,0) {loginForm: FormGroup};
  \node[cp, anchor=west] (cu) at (2.9,-1) {user: FormControl};
  \node[cp, anchor=west] (cpw) at (2.9,-2) {pass: FormControl};
  \draw[l] (f.east) -- node[mcod, above]{[formGroup]} (g.west);
  \draw[l] (iu.east) -- node[mcod, above]{formControlName} (cu.west);
  \draw[l] (ip.east) -- node[mcod, above]{formControlName} (cpw.west);
\end{tikzpicture}
\end{immagine}
\vspace{6pt}
\renewcommand{\wtdialogohost}{localhost:4200}
\begin{immagine}[Premendo \textit{Submit} viene invocato \code{onSubmit()}.]
\wtdialogo{Username: alice\\\hspace*{3em}Password: 12345}{\wtbottone[blue!70!black][white][blue!70!black]{OK}}
\end{immagine}
\end{affianco}

Qui la struttura del form è dichiarata esplicitamente nella classe: è più
verboso, ma rende il form ispezionabile e testabile dal codice, e permette di
aggiungere validatori in modo strutturato.

\section{Servizi e iniezione delle dipendenze}

Sappiamo che le classi del nostro software dovrebbero fare \textbf{una cosa
sola} e farla bene: aumenta modularità e riusabilità. I componenti Angular
dovrebbero occuparsi soltanto dell'esperienza utente: template, stili e la
logica che fa da mediatore fra la vista e la logica applicativa.

\dfn{Servizio}{
  I \term{servizi} sono classi che forniscono funzionalità utilizzabili in più
  parti dell'applicazione. Non sono associati direttamente all'esperienza
  utente: niente template, niente stili, solo logica invocabile all'occorrenza.}

Sono ideali per compiti come recuperare dati da una API REST, registrare log o
validare dati: attività che con ogni probabilità serviranno in punti diversi
dell'applicazione.

\dfn{Dependency Injection}{
  Angular sfrutta il pattern della \term{iniezione delle dipendenze} per
  permettere al codice di ottenere facilmente istanze delle proprie dipendenze.
  L'obiettivo è \textbf{separare} la responsabilità di costruire le dipendenze
  da quella di usarle, riducendo l'accoppiamento.}

Nel sistema di iniezione di Angular esistono due ruoli: il \textbf{consumatore}
di dipendenze (il codice che scriviamo noi) e il \textbf{fornitore} di
dipendenze (gestito generalmente da Angular).

\subsection{Creare un servizio}

\begin{lstlisting}[style=shell]
ng generate service calculator
\end{lstlisting}

\begin{affianco}[0.46][c]
\codicepiccolo
\begin{lstlisting}[style=ts, name=calculator.service.ts]
import { Injectable } from '@angular/core';

@Injectable({
  providedIn: 'root'  // disponibile ovunque
})
export class CalculatorService {
  sum(x: number, y: number): number { return x + y; }
}
\end{lstlisting}
\risultato
\begin{immagine}[Una sola istanza, creata da Angular\\
                 e consegnata a chi la chiede.]
\begin{tikzpicture}[
    a/.style={-{Stealth[length=4pt,width=3.5pt]}, line width=0.75pt,
      draw=neutral!45!black}]
  \node[mfun] (s) at (2.4,0) {CalculatorService};
  \node[fit=(s), draw=primary!45!white, fill=primary!5!white, rounded corners=5pt,
        inner xsep=10pt, inner ysep=12pt, yshift=4pt] (inj) {};
  \node[mscopolab, anchor=north west] at (inj.north west) {iniettore radice};
  \node[mfun] (s) at (2.4,0) {CalculatorService};
  \node[wtnodeplain, font=\FiraCodeNL\scriptsize] (c1) at (0,-1.9) {HomePageComponent};
  \node[wtnodeplain, font=\FiraCodeNL\scriptsize] (c2) at (4.8,-1.9) {AltroComponent};
  \draw[a] (inj.south) ++(-0.3,0) -- node[mcod, left=2pt]{inject()} (c1.north);
  \draw[a] (inj.south) ++(0.3,0) -- node[mcod, right=2pt]{inject()} (c2.north);
\end{tikzpicture}
\end{immagine}
\end{affianco}

Il decoratore \code{@Injectable} marca la classe come qualcosa che può essere
iniettato in un componente come dipendenza.

\subsection{Iniettare un servizio}

Ci sono due modi. Con la funzione \code{inject()}:

\begin{affianco}[0.6][c]
\codicepiccolo
\begin{lstlisting}[style=ts, name=home-page.component.ts]
import { Component, inject } from '@angular/core';
import { CalculatorService } from '../calculator.service';

@Component({
  selector: 'app-home-page',
  standalone: true,
  template: "<p>The sum is <strong>{{sum}}</strong></p>",
})
export class HomePageComponent {
  private calculatorService = inject(CalculatorService);
  sum = this.calculatorService.sum(1234, 4321);
}
\end{lstlisting}
\risultato
\begin{browser}[localhost:4200]
  The sum is \textbf{5555}
\end{browser}
\wtnota{$1234 + 4321 = 5555$}
\end{affianco}

Oppure dichiarandolo nel costruttore della classe:

\begin{lstlisting}[style=ts, name=home-page.component.ts]
export class HomePageComponent {
  constructor(private calculatorService: CalculatorService) {}
  sum = this.calculatorService.sum(1234, 4321);
}
\end{lstlisting}

In nessuno dei due casi il componente costruisce il servizio con \code{new}: è
precisamente questo il punto dell'iniezione delle dipendenze, e ciò che rende il
componente facilmente testabile sostituendo il servizio con una versione finta.

\section{Proteggere le rotte}

Un compito frequente è assicurarsi che solo gli utenti autorizzati accedano a
certe rotte. Con il Router di Angular si usano le \term{Route Guard}: funzioni
invocate da Angular per determinare se l'utente possa compiere certe operazioni
su una rotta (\code{canActivate}, \code{canMatch}, \code{canLoad},
\code{canDeactivate}).

Le guardie si indicano come proprietà degli oggetti \code{Route}. Se ne possono
specificare più di una: Angular le esegue \textbf{in parallelo}, e l'azione è
consentita solo se tutte restituiscono \code{true}.

\begin{lstlisting}[style=shell]
ng generate guard authorization
\end{lstlisting}

\begin{lstlisting}[style=ts, name=authorization.guard.ts]
import { inject } from '@angular/core';
import { CanActivateFn, Router } from '@angular/router';
import { AuthService } from './auth.service';

export const authorizationGuard: CanActivateFn = (route, state) => {
  const router = inject(Router);
  const authService = inject(AuthService);

  if (authService.isUserLoggedIn()) {
    return true;                      // utente autorizzato
  } else {
    router.navigateByUrl("/login");   // reindirizza al login
    return false;
  }
};
\end{lstlisting}

\begin{affianco}[0.4][c]
\begin{lstlisting}[style=ts, name=app.routes.ts]
{
  path: "counter",
  title: "Counter",
  component: CounterPageComponent,
  canActivate: [authorizationGuard]
}
\end{lstlisting}
\risultato
\begin{immagine}
\begin{tikzpicture}[
    dec/.style={diamond, aspect=2.4, draw=primary!60!black, fill=primary!8!white,
      line width=0.7pt, font=\Lato\scriptsize, inner sep=1pt, align=center},
    esito/.style={rounded corners=3pt, draw=neutral!50!black, fill=white,
      line width=0.7pt, font=\FiraCodeNL\scriptsize, inner sep=3pt, align=center},
    a/.style={-{Stealth[length=4pt,width=3.5pt]}, line width=0.75pt,
      draw=neutral!45!black}]
  \node[esito] (n) at (0,0) {navigazione verso /counter};
  \node[dec] (q) at (0,-1.4) {utente\\autenticato?};
  \node[esito, draw=mok] (ok) at (0,-2.95) {CounterPageComponent};
  \node[esito, draw=merr, anchor=west] (no) at (2.9,-1.4) {navigateByUrl("/login")};
  \draw[a] (n) -- node[mnota, right]{\code{authorizationGuard}} (q);
  \draw[a] (q) -- node[mnota, right]{sì: \code{true}} (ok);
  \draw[a] (q) -- node[mnota, above]{no: \code{false}} (no);
\end{tikzpicture}
\end{immagine}
\end{affianco}

Una guardia può restituire un valore booleano oppure un oggetto \code{UrlTree}:

\begin{itemize}
  \item \code{true}: la guardia è soddisfatta, la navigazione può proseguire;
  \item \code{false}: la guardia non è soddisfatta, l'azione non è consentita;
  \item un \code{UrlTree}: il Router deve reindirizzare a quella URL.
\end{itemize}

\warning{Con più guardie, restituire un UrlTree}{
  Quando si impostano più guardie bisogna sempre restituire un \code{UrlTree}
  per reindirizzare, e \textbf{non} invocare direttamente
  \code{Router.navigate()} come nell'esempio precedente. Con più guardie
  eseguite in parallelo, se due di esse reindirizzano si possono avere
  \textit{race condition}.}

\error{Le guardie non sono sicurezza}{
  Una Route Guard gira \textbf{nel browser}, cioè su una macchina interamente
  sotto il controllo dell'utente. Serve a migliorare l'esperienza (non mostrare
  pagine inutilizzabili), non a proteggere i dati. La vera autorizzazione va
  sempre imposta \textbf{sul server}, come abbiamo fatto nel capitolo sulle API
  REST.}

\section{HttpClient}

Angular fornisce una API per effettuare richieste HTTP, implementata nella
classe di servizio \code{HttpClient}. Va configurata tramite iniezione delle
dipendenze.

\begin{lstlisting}[style=ts, name=app.config.ts]
export const appConfig: ApplicationConfig = {
  providers: [
    provideHttpClient(),
    provideRouter(routes)
  ]
};
\end{lstlisting}

\begin{lstlisting}[style=ts, name=rest-backend.service.ts]
import { HttpClient } from '@angular/common/http';
import { Injectable } from '@angular/core';

@Injectable({ providedIn: 'root' })
export class RestBackendService {
  constructor(private http: HttpClient) {}
  // il servizio puo' ora effettuare richieste tramite `this.http`
}
\end{lstlisting}

\begin{esempio}{Una richiesta HTTP}
\begin{affianco}[0.5][c]
\codicepiccolo
\begin{lstlisting}[style=ts, name=cat-facts.component.ts]
@Component({
  selector: 'app-cat-facts',
  standalone: true,
  template: `
    <h1>Cat Facts</h1>
    <button (click)="loadCatFact()">
      Click here to load a cat fact
    </button>
    <p>{{catFact}}</p>
  `
})
export class CatFactsComponent {
  constructor(private http: HttpClient) {}

  catFact: string = "";

  loadCatFact() {
    let url = "https://cat-fact.herokuapp.com/facts/random";
    this.http.get<{ text: string }>(url).subscribe(data => {
      this.catFact = data.text;
    });
  }
}
\end{lstlisting}
\risultato
\begin{browser}[localhost:4200]
  {\Lora\bfseries\LARGE Cat Facts\par}\vspace{6pt}
  \wtbottone{Click here to load a cat fact}\par\vspace{6pt}
  Cats sleep for around 13 to 16 hours a day.
\end{browser}
\begin{immagine}
\begin{tikzpicture}[
    a/.style={-{Stealth[length=4pt,width=3.5pt]}, line width=0.75pt,
      draw=neutral!45!black}]
  \node[wtnodeplain, font=\Lato\scriptsize] (c) at (0,0) {clic};
  \node[wtnodeplain, font=\FiraCodeNL\scriptsize] (g) at (2.6,0) {http.get(url)};
  \node[wtnodeplain, font=\FiraCodeNL\scriptsize] (s) at (2.6,-1.1) {subscribe(data => \ldots)};
  \node[wtnodeok, font=\FiraCodeNL\scriptsize, inner sep=3pt] (v) at (2.6,-2.2)
    {<p>\{\{catFact\}\}</p>};
  \draw[a] (c) -- (g);
  \draw[a] (g) -- node[mnota, right]{risposta JSON} (s);
  \draw[a] (s) -- node[mnota, right]{\code{this.catFact = data.text}} (v);
\end{tikzpicture}
\end{immagine}
\end{affianco}
\end{esempio}

\section{Observable}

\code{HttpClient} fornisce metodi corrispondenti ai vari verbi HTTP, e ciascuno
restituisce un \term{Observable} della libreria RxJS (\textit{Reactive
Extensions Library for JavaScript}). Angular usa gli Observable per gestire
molte operazioni asincrone.

\dfn{Observable}{
  Un \term{Observable} è per certi versi simile a una Promise, con una
  differenza essenziale: una \textbf{Promise} può risolversi (o essere
  rifiutata) \textbf{una volta sola}, cioè emette in modo asincrono un
  \textbf{singolo} valore o un errore; un \textbf{Observable} può emettere in
  modo asincrono \textbf{più valori nel tempo}.}

\begin{center}
\renewcommand{\arraystretch}{1.4}
\begin{tabular}{@{}>{\raggedright\arraybackslash}p{6.3cm}@{\hspace{1.0cm}}>{\raggedright\arraybackslash}p{6.3cm}@{}}
  \textbf{Promise} & \textbf{Observable} \\
  Un solo valore, una sola volta. & Zero, uno o molti valori nel tempo. \\
  Si consuma con \code{.then()} / \code{.catch()}. & Si consuma
    \term{sottoscrivendosi} con \code{.subscribe()}. \\
  Parte subito, alla creazione. & Parte alla sottoscrizione. \\
\end{tabular}
\end{center}

\subsection{Sottoscriversi a un Observable}

Come con le Promise, bisogna fornire le callback da eseguire al verificarsi
degli eventi. Lo si fa \textbf{sottoscrivendosi}.

\begin{affianco}[0.46][c]
\codicepiccolo
\begin{lstlisting}[style=ts, name=subscribe.ts]
observable.subscribe({
  next(value) {   // a ogni valore emesso
    console.log(`Observable emitted value ${value}`);
  },
  error(err) {    // se emette un errore
    console.log(`Observable emitted an Error: ${err}`);
  },
  complete() {    // quando ha finito
    console.log("Observable finished emitting values");
  }
});
\end{lstlisting}
\risultato
\begin{immagine}[Ogni valore emesso invoca \code{next}; poi l'Observable\\
                 termina con \code{complete} oppure con \code{error}.]
\begin{tikzpicture}[
    m/.style={circle, draw=primary!65!black, fill=primary!14!white,
      line width=0.7pt, minimum size=0.42cm, inner sep=0pt,
      font=\FiraCodeNL\tiny},
    t/.style={-{Stealth[length=5pt,width=4pt]}, line width=0.8pt,
      draw=neutral!50!black},
    fine/.style={line width=1.4pt, draw=mok},
    err/.style={font=\Lato\normalsize\bfseries, text=merr, inner sep=0pt}]
  \draw[t] (0,0) -- (6,0) node[mnota, right]{tempo};
  \foreach \x/\v in {0.8/3,2.1/7,3.4/1} \node[m] (v\x) at (\x,0) {\v};
  \draw[fine] (4.6,-0.25) -- (4.6,0.25);
  \node[mcod, anchor=south] at (2.1,0.3) {next(v)};
  \node[mcod, anchor=south] at (4.6,0.3) {complete()};
  \draw[t] (0,-1.2) -- (6,-1.2);
  \foreach \x/\v in {0.8/4,2.1/9} \node[m] at (\x,-1.2) {\v};
  \node[err] at (3.4,-1.2) {$\times$};
  \node[mcod, anchor=north] at (3.4,-1.45) {error(err)};
\end{tikzpicture}
\end{immagine}
\end{affianco}

\subsection{Creare un Observable}

L'argomento passato al costruttore è la funzione da eseguire al momento della
sottoscrizione.

\begin{affianco}[0.52][c]
\codicepiccolo
\begin{lstlisting}[style=ts, name=observables_example.ts]
import { Observable } from "rxjs";

const observable = new Observable<number>((subscriber) => {
  setTimeout(() => {
    let num = Math.random();
    if (num < 0.5) {
      subscriber.next(num);
    } else {
      subscriber.error(new Error(`You were unlucky! ${num}`));
    }
    subscriber.complete();
  }, 2000);
});
\end{lstlisting}
\risultato
\begin{immagine}[I due esiti possibili, due secondi\\
                 dopo la sottoscrizione.]
\begin{tikzpicture}[
    m/.style={circle, draw=primary!65!black, fill=primary!14!white,
      line width=0.7pt, minimum size=0.42cm, inner sep=0pt,
      font=\FiraCodeNL\tiny},
    t/.style={-{Stealth[length=5pt,width=4pt]}, line width=0.8pt,
      draw=neutral!50!black},
    fine/.style={line width=1.4pt, draw=mok},
    err/.style={font=\Lato\normalsize\bfseries, text=merr, inner sep=0pt}]
  \node[mnota, anchor=east] at (-0.1,0) {\code{num < 0.5}};
  \draw[t] (0,0) -- (4.2,0);
  \node[m, minimum size=0.5cm] at (2.6,0) {num};
  \draw[fine] (3.15,-0.25) -- (3.15,0.25);
  \node[mnota, anchor=east] at (-0.1,-1.1) {altrimenti};
  \draw[t] (0,-1.1) -- (4.2,-1.1);
  \node[err] at (2.6,-1.1) {$\times$};
  \draw[wtdash] (2.6,0.45) -- (2.6,-1.5) node[mnota, below]{2 s};
\end{tikzpicture}
\end{immagine}
\end{affianco}

\begin{console}[Terminale]
$ node observables_example.js
Observable emitted value 0.33397894647355697
Observable finished emitting values

$ node observables_example.js
Observable emitted an Error: Error: You were unlucky! 0.620240538767127
\end{console}

\warning{Un errore è una conclusione}{
  Nella seconda esecuzione \code{complete()} \textbf{non} viene invocata: un
  errore termina l'Observable, e le due callback si escludono a vicenda. È lo
  stesso comportamento di un'eccezione che interrompe un flusso.}

\subsection{Emettere più valori}

\begin{affianco}[0.6][c]
\codicepiccolo
\begin{lstlisting}[style=ts, name=multiple_values.ts]
import { Observable, Subscriber } from "rxjs";

const observable = new Observable<number>(subscriber => {
  helper(subscriber, 1);
});

function helper(subscriber: Subscriber<number>, level: number) {
  setTimeout(() => {
    let num = Math.random();
    if (level > 3) { subscriber.complete(); return; }

    if (num < 0.8) {
      subscriber.next(num);
      helper(subscriber, level + 1);   // ricorsione
    } else {
      subscriber.error(new Error(`You were unlucky! ${num}`));
    }
  }, 1000);
}
\end{lstlisting}
\risultato
\begin{immagine}[Un valore al secondo; al quarto\\
                 livello la sequenza termina.]
\begin{tikzpicture}[
    m/.style={circle, draw=primary!65!black, fill=primary!14!white,
      line width=0.7pt, minimum size=0.42cm, inner sep=0pt,
      font=\FiraCodeNL\tiny},
    t/.style={-{Stealth[length=5pt,width=4pt]}, line width=0.8pt,
      draw=neutral!50!black},
    fine/.style={line width=1.4pt, draw=mok},
    err/.style={font=\Lato\normalsize\bfseries, text=merr, inner sep=0pt}]
  \draw[t] (0,0) -- (4.6,0);
  \foreach \x/\l in {1/1,2/2,3/3} {
    \node[m] at (\x,0) {};
    \node[mindice, anchor=north] at (\x,-0.3) {\x\,s};
    \node[mindice, anchor=south] at (\x,0.3) {level \l};}
  \draw[fine] (4,-0.25) -- (4,0.25);
  \node[mindice, anchor=north] at (4,-0.3) {4\,s};
\end{tikzpicture}
\end{immagine}
\end{affianco}

\begin{console}[Terminale]
$ node multiple_values.js
Observable emitted value 0.6982903488018062
Observable emitted value 0.5713928193844571
Observable emitted value 0.005029574874092724
Observable finished emitting values
\end{console}

\section{Interceptor}

\code{HttpClient} si può configurare in molti modi: per esempio con
\code{withFetch()} per usare la \code{fetch} API al posto di
\code{XMLHttpRequest}, oppure aggiungendo \term{interceptor}.

\begin{affianco}[0.52][c]
\codicepiccolo
\begin{lstlisting}[style=ts, name=app.config.ts]
export const appConfig: ApplicationConfig = {
  providers: [
    provideHttpClient(
      withFetch(),
      withInterceptors([loggingInterceptor, authInterceptor])
    ),
    provideRouter(routes)
  ]
};
\end{lstlisting}
\risultato
\begin{immagine}
\begin{tikzpicture}[
    a/.style={-{Stealth[length=4pt,width=3.5pt]}, line width=0.75pt,
      draw=neutral!45!black},
    n/.style={wtnodeplain, font=\FiraCodeNL\scriptsize, minimum width=3.4cm}]
  \node[n, font=\Lato\scriptsize] (c) at (0,0) {richiesta del componente};
  \node[wtnodeacc, font=\FiraCodeNL\scriptsize, minimum width=3.4cm, inner sep=3pt]
    (l) at (0,-1) {loggingInterceptor};
  \node[wtnodeacc, font=\FiraCodeNL\scriptsize, minimum width=3.4cm, inner sep=3pt]
    (au) at (0,-2) {authInterceptor};
  \node[n] (f) at (0,-3) {fetch $\to$ server};
  \draw[a] (c) -- (l);
  \draw[a] (l) -- node[mcod, right]{next(req)} (au);
  \draw[a] (au) -- node[mcod, right]{next(req)} (f);
\end{tikzpicture}
\end{immagine}
\end{affianco}

Gli interceptor sono una forma di \textbf{middleware} supportata da
\code{HttpClient}: stesso concetto dei middleware di Express visti nel capitolo
13, ma applicato alle richieste \textbf{in uscita}. Sono funzioni che ricevono
la richiesta corrente e una funzione \code{next}, che rappresenta il passo
successivo della catena.

\begin{affianco}[0.62][c]
\codicepiccolo
\begin{lstlisting}[style=ts, name=interceptors.ts]
function log(req: HttpRequest<unknown>, next: HttpHandlerFn) {
  console.log(req.url);
  return next(req);
}
\end{lstlisting}
\risultato
\begin{console}
http://localhost:3000/todos
\end{console}
\end{affianco}

Casi d'uso tipici:

\begin{itemize}
  \item aggiungere header di autenticazione alle richieste (per esempio i token
        JWT del capitolo 15);
  \item mettere in cache le risposte per un certo tempo;
  \item registrare le richieste in uscita.
\end{itemize}

\section{Change detection}

Quando una proprietà di un componente cambia, il suo template viene aggiornato
automaticamente da Angular. Aggiornare i template \textbf{solo quando serve} e
farlo in modo efficiente (aggiornando solo le parti necessarie) è essenziale per
le prestazioni di una SPA.

Come fa Angular a sapere quando deve aggiornare la vista?

In sostanza, quando accade qualcosa nell'applicazione (eventi del DOM,
operazioni asincrone) si attiva il meccanismo di \term{change detection}.
Angular percorre tutti i componenti della gerarchia e, per ciascuno, verifica se
il suo stato è cambiato e se il cambiamento influisce sulla vista; in tal caso
aggiorna la porzione di DOM corrispondente al template del componente.

\section{Signal}

I \term{signal} sono stati introdotti in Angular 17. Tracciano in modo
\textbf{granulare} come e dove lo stato viene usato nell'applicazione,
permettendo ad Angular di ottimizzare ulteriormente gli aggiornamenti.

\dfn{Signal}{
  Si può pensare a un \term{signal} come a una variabile capace anche di
  \textbf{notificare} i consumatori interessati quando il proprio valore cambia.
  È un modo di realizzare la \term{programmazione reattiva}, una forma di
  programmazione dichiarativa basata sull'elaborazione asincrona degli eventi.}

Usare i signal permette ad Angular di evitare controlli di change detection non
necessari e di aggiornare solo le parti del DOM effettivamente cambiate.

\subsection{Creare e aggiornare}

I signal possono essere scrivibili o in sola lettura. Quelli scrivibili si
creano con la funzione \code{signal()}, passando un valore iniziale, e offrono
\code{.set()} per impostare un nuovo valore e \code{.update()} per calcolarlo a
partire dal precedente.

\begin{affianco}[0.52][c]
\codicepiccolo
\begin{lstlisting}[style=ts, name=counter-signal.component.ts]
import { Component, signal } from '@angular/core';

@Component({
  selector: 'app-counter-signal',
  standalone: true,
  template: `
    <h1>Counter with Signals</h1>
    <button (click)="handleClick()">Click to Increment</button>
    <button (click)="handleReset()">Click to Reset</button>
    <p>{{count()}}</p>   <!-- si notino le parentesi! -->
  `
})
export class CounterSignalComponent {
  count = signal(0);

  handleClick() { this.count.update(value => value + 1); }
  handleReset() { this.count.set(0); }
}
\end{lstlisting}
\risultato
\begin{browser}[localhost:4200]
  {\Lora\bfseries\LARGE Counter with Signals\par}\vspace{6pt}
  \wtbottone{Click to Increment}\,\wtbottone{Click to Reset}\par\vspace{6pt}
  3
\end{browser}
\begin{immagine}
\begin{tikzpicture}[
    a/.style={-{Stealth[length=4pt,width=3.5pt]}, line width=0.75pt,
      draw=neutral!45!black}]
  \node[mvar, minimum width=1.3cm] (s) at (0,0) {3};
  \node[mnome, anchor=south] at (s.north) {count};
  \node[mcod, anchor=east] (u) at (-1.6,0.45) {update(v => v + 1)};
  \node[mcod, anchor=east] (r) at (-1.6,-0.45) {set(0)};
  \node[mcod, anchor=west] (t) at (1.6,0) {\{\{count()\}\}};
  \draw[a] (u.east) -- (s.west); \draw[a] (r.east) -- (s.west);
  \draw[a] (s.east) -- node[mnota, above]{legge} (t.west);
\end{tikzpicture}
\end{immagine}
\end{affianco}

Nel template il signal si legge \textbf{invocandolo}: \code{count()}, non
\code{count}.

\subsection{Signal calcolati}

I \term{signal calcolati} sono in sola lettura e derivano il proprio valore da
quello di altri signal. Si definiscono con la funzione \code{computed()} e una
funzione di derivazione; Angular li mantiene aggiornati automaticamente, in modo
reattivo. Invocare \code{.set()} o \code{.update()} su di essi produce un
errore.

\begin{affianco}[0.48][c]
\begin{lstlisting}[style=ts]
export class CounterSignalComponent {
  count = signal(0);
  double = computed(() => this.count() * 2);
}
\end{lstlisting}

\begin{lstlisting}[style=html]
<p>Count: {{count()}} - Double: {{double()}}</p>
\end{lstlisting}
\risultato
\begin{immagine}
\begin{tikzpicture}[
    a/.style={-{Stealth[length=4pt,width=3.5pt]}, line width=0.75pt,
      draw=neutral!45!black}]
  \node[mvar, minimum width=1.2cm] (c) at (0,0) {3};
  \node[mnome, anchor=south] at (c.north) {count};
  \node[mvar, minimum width=1.2cm, fill=primary!6!white] (d) at (4.6,0) {6};
  \node[mnome, anchor=south] at (d.north) {double};
  \draw[a] (c) -- node[mcod, below]{() => count() * 2} (d);
  \node[mnota, anchor=north] at ([yshift=-3pt]d.south) {sola lettura};
\end{tikzpicture}
\end{immagine}
\begin{browser}[localhost:4200]
  Count: 3 - Double: 6
\end{browser}
\end{affianco}

\subsection{Effetti}

Un \term{effetto} è una funzione che viene eseguita ogni volta che certi signal
cambiano. Si crea con la funzione \code{effect()}.

\begin{affianco}[0.56][c]
\codicepiccolo
\begin{lstlisting}[style=ts, name=counter-signal.component.ts]
export class CounterSignalComponent {
  count = signal(0);
  double = computed(() => this.count() * 2);

  constructor() {
    effect(() => {
      console.log(`Count value changed: ${this.count()}`);
    });
  }

  handleClick() { this.count.update(value => value + 1); }
  handleReset() { this.count.set(0); }
}
\end{lstlisting}
\risultato
\begin{immagine}[Angular registra i signal letti dall'effetto\\
                 e lo riesegue quando cambiano.]
\begin{tikzpicture}[
    a/.style={-{Stealth[length=4pt,width=3.5pt]}, line width=0.75pt,
      draw=neutral!45!black}]
  \node[mvar, minimum width=1.2cm] (c) at (0,0) {2};
  \node[mnome, anchor=south] at (c.north) {count};
  \node[mfun] (e) at (2.6,0) {effect(\ldots)};
  \draw[a] (c) -- (e);
\end{tikzpicture}
\end{immagine}
\begin{console}
Count value changed: 0
Count value changed: 1
Count value changed: 2
\end{console}
\end{affianco}

Gli effetti vengono eseguiti \textbf{almeno una volta}: è per questo che la
prima riga compare senza che nessuno abbia premuto il pulsante. Quando vengono
eseguiti, tengono traccia di tutti i signal che leggono, e ogni volta che uno di
questi cambia l'effetto viene rieseguito.

\section{Mettere tutto insieme: la lista come SPA}

Chiudiamo il percorso riscrivendo la nostra applicazione come Single Page App in
Angular, appoggiandosi al \textbf{backend REST} sviluppato nel capitolo 15. È
l'occasione per vedere come i pezzi studiati finora si incastrano.

\subsection{Le rotte}

\begin{affianco}[0.46][c]
\codicepiccolo
\begin{lstlisting}[style=ts, name=app.routes.ts]
export const routes: Routes = [
  {
    path: "home", component: HomepageComponent,
    title: "To-do List Angular App"
  }, {
    path: "login", component: LoginComponent,
    title: "Login | To-do List Angular App"
  }, {
    path: "todos", component: TodoPageComponent,
    title: "To-do List | To-do List Angular App",
    canActivate: [authGuard]  // protetta
  }, {
    path: "todos/:id", component: TodoDetailComponent,
    title: "To-do Detail | To-do List Angular App",
    canActivate: [authGuard]  // protetta e parametrica
  }
];
\end{lstlisting}
\risultato
\begin{immagine}
\begin{tikzpicture}[
    u/.style={mcod, draw=neutral!55!black, fill=white, line width=0.7pt,
      minimum width=1.8cm, minimum height=0.5cm},
    a/.style={-{Stealth[length=4pt,width=3.5pt]}, line width=0.75pt,
      draw=neutral!45!black}]
  \foreach \p/\c/\g [count=\i] in {{/home}/HomepageComponent/0,
      {/login}/LoginComponent/0, {/todos}/TodoPageComponent/1,
      {/todos/:id}/TodoDetailComponent/1} {
    \node[u] (u\i) at (0,-0.85*\i) {\p};
    \node[mfun, anchor=west] (c\i) at (3.0,-0.85*\i) {\c};
    \draw[a] (u\i) -- (c\i);
    \ifnum\g=1
      \node[mcod, text=merr, fill=white, inner sep=1pt] at (2.2,-0.85*\i)
        {authGuard};
    \fi}
\end{tikzpicture}
\end{immagine}
\end{affianco}

\subsection{La guardia}

\begin{affianco}[0.54][c]
\codicepiccolo
\begin{lstlisting}[style=ts, name=auth.guard.ts]
export const authGuard: CanActivateFn = (route, state) => {
  const authService = inject(AuthService);
  const toastr = inject(ToastrService);
  const router = inject(Router);

  if (authService.isUserAuthenticated()) {
    return true;
  } else {
    toastr.warning("Effettua l'accesso per usare questa funzione",
                   "Non autorizzato!");
    return router.parseUrl("/login");     // restituisce un UrlTree
  }
};
\end{lstlisting}
\risultato
\begin{immagine}[La notifica di \code{toastr.warning}; intanto\\
                 il Router passa a \code{/login}.]
\wttoast[warning]{Non autorizzato!}{Effettua l'accesso per usare questa funzione}
\end{immagine}
\end{affianco}

Si noti \code{router.parseUrl()}: è la forma corretta di reindirizzamento
discussa poco fa, quella che evita le \textit{race condition} con più guardie.

\subsection{Il componente di login}

\begin{lstlisting}[style=ts, name=login.component.ts]
export class LoginComponent {
  toastr      = inject(ToastrService);
  router      = inject(Router);
  restService = inject(RestBackendService);
  authService = inject(AuthService);

  submitted = false;

  loginForm = new FormGroup({
    user: new FormControl('', [Validators.required]),
    pass: new FormControl('', [
      Validators.required,
      Validators.minLength(4),
      Validators.maxLength(16)
    ])
  });

  handleLogin() {
    this.submitted = true;

    if (this.loginForm.invalid) {
      this.toastr.error("Dati non validi!", "Attenzione");
      return;
    }

    this.restService.login({
      usr: this.loginForm.value.user as string,
      pwd: this.loginForm.value.pass as string,
    }).subscribe({
      next:  (token) => { this.authService.updateToken(token as string); },
      error: (err)   => { this.toastr.error("Credenziali non valide"); },
      complete: () => {
        this.toastr.success(`Benvenuto ${this.loginForm.value.user}!`);
        this.router.navigateByUrl("/todos");
      }
    });
  }
}
\end{lstlisting}

Nel template, gli errori di validazione si mostrano solo quando servono, con la
sintassi di controllo di flusso di Angular 17:

\begin{affianco}[0.5][c]
\codicepiccolo
\begin{lstlisting}[style=html, name=login.component.html]
<label for="pass">Password</label>
<input type="password" formControlName="pass" id="pass" required>

@if (submitted && loginForm.controls.pass.errors) {
  @if (loginForm.controls.pass.errors['required']) {
    <p class="form-error">
      La password e' obbligatoria.
    </p>
  }
  @if (loginForm.controls.pass.errors['minlength']) {
    <p class="form-error">
      La password deve contenere almeno 4 caratteri.
    </p>
  }
}
\end{lstlisting}
\risultato
\begin{browser}[localhost:4200/login]
  Password\par\vspace{2pt}
  \wtcampo[3cm]{$\bullet\bullet\bullet$}\par\vspace{4pt}
  {\color{css-red}La password deve contenere almeno 4 caratteri.\par}
\end{browser}
\begin{immagine}[Inviato con una password di tre caratteri.]
\wttoast[error]{Attenzione}{Dati non validi!}
\end{immagine}
\end{affianco}

\subsection{Il servizio di autenticazione, con i signal}

\begin{lstlisting}[style=ts, name=auth.service.ts]
type AuthState = {
  user: string | null,
  token: string | null,
  isAuthenticated: boolean
}

@Injectable({ providedIn: 'root' })
export class AuthService {

  authState: WritableSignal<AuthState> = signal<AuthState>({
    user:  this.getUser(),
    token: this.getToken(),                        // dal localStorage, se c'e'
    isAuthenticated: this.verifyToken(this.getToken())   // non scaduto?
  });

  user            = computed(() => this.authState().user);
  token           = computed(() => this.authState().token);
  isAuthenticated = computed(() => this.authState().isAuthenticated);

  constructor() {
    // un effetto tiene allineati i dati con il localStorage
    effect(() => {
      const token = this.authState().token;
      const user  = this.authState().user;

      if (token !== null) localStorage.setItem("token", token);
      else                localStorage.removeItem("token");

      if (user !== null)  localStorage.setItem("user", user);
      else                localStorage.removeItem("user");
    });
  }
}
\end{lstlisting}

\begin{affianco}[0.42][c]
Qui i signal danno il meglio: lo stato di autenticazione è un solo signal
scrivibile, da cui derivano tre signal calcolati, e un \textbf{effetto} tiene
sincronizzato il \code{localStorage} senza che nessuno debba ricordarsi di
chiamarlo.
\risultato
\begin{immagine}
\begin{tikzpicture}[
    a/.style={-{Stealth[length=4pt,width=3.5pt]}, line width=0.75pt,
      draw=neutral!45!black},
    cs/.style={mvar, minimum width=2.6cm, fill=primary!6!white}]
  \node[mvar, minimum width=2.4cm, line width=1pt] (s) at (0,0) {authState};
  \node[mnota, anchor=north] at ([yshift=-1pt]s.south) {scrivibile};
  \node[cs] (u) at (4,0.9) {user};
  \node[cs] (t) at (4,0) {token};
  \node[cs] (i) at (4,-0.9) {isAuthenticated};
  \foreach \x in {u,t,i} \draw[a] (s.east) -- (\x.west);
  \node[mnota, anchor=south] at (u.north) {\code{computed}};
  \node[mfun] (e) at (0,-1.9) {effect(\ldots)};
  \node[wtnodeplain, font=\FiraCodeNL\scriptsize] (ls) at (4,-1.9) {localStorage};
  \draw[a] ([yshift=-8pt]s.south) -- (e.north);
  \draw[a] (e) -- node[mnota, above]{sincronizza} (ls);
\end{tikzpicture}
\end{immagine}
\end{affianco}

\subsection{Il servizio REST e l'interceptor}

\begin{lstlisting}[style=ts, name=rest-backend.service.ts]
export interface AuthRequest {
  usr: string,
  pwd: string
}

@Injectable({ providedIn: 'root' })
export class RestBackendService {
  url = "http://localhost:3000";

  constructor(private http: HttpClient) {}

  httpOptions = {
    headers: new HttpHeaders({ 'Content-Type': 'application/json' })
  };

  login(loginRequest: AuthRequest) {
    const url = `${this.url}/auth`;
    return this.http.post<string>(url, loginRequest, this.httpOptions);
  }
}
\end{lstlisting}

\begin{affianco}[0.5][c]
\codicepiccolo
\begin{lstlisting}[style=ts, name=auth.interceptor.ts]
function authInterceptor(request: HttpRequest<unknown>,
                         next: HttpHandlerFn) {
  const authService = inject(AuthService);
  const token = authService.getToken();

  if (token) {
    // una copia della richiesta, con il token
    request = request.clone({
      setHeaders: {
        Authorization: 'Bearer ' + token
      }
    });
  }

  return next(request);
}
\end{lstlisting}
\risultato
\begin{immagine}[Le richieste sono immutabili: si modifica una copia.]
\begin{tikzpicture}[
    a/.style={-{Stealth[length=4pt,width=3.5pt]}, line width=0.75pt,
      draw=neutral!45!black},
    rq/.style={draw=neutral!50!black, fill=white, line width=0.7pt,
      rounded corners=2pt, inner sep=4pt, align=left, font=\FiraCodeNL\scriptsize}]
  \node[rq] (o) at (0,0) {GET /todos\\Content-Type: application/json};
  \node[mnota, anchor=south west] at (o.north west) {\textbf{originale}};
  \node[rq] (c) at (0,-2.1) {GET /todos\\Content-Type: application/json\\
    \textcolor{accent!75!black}{Authorization: Bearer eyJhbGci\ldots}};
  \node[mnota, anchor=south west] at (c.north west) {\textbf{copia}};
  \draw[a] (o.south) -- node[mcod, right]{clone(\{ setHeaders \})} ([yshift=10pt]c.north);
\end{tikzpicture}
\end{immagine}
\end{affianco}

L'interceptor è il pezzo che chiude il cerchio con il capitolo 15: il token JWT
viene allegato \textbf{automaticamente} a ogni richiesta uscente, e nessun
componente deve occuparsene.

\subsection{La lista e i suoi elementi}

\begin{lstlisting}[style=ts, name=todo-page.component.ts]
export class TodoPageComponent {
  createTodoSubmitted = false;

  restService = inject(RestBackendService);
  toastr      = inject(ToastrService);
  router      = inject(Router);

  todos: TodoItem[] = [];

  newTodoForm = new FormGroup({
    todo: new FormControl('', [Validators.required])
  });

  ngOnInit() { this.fetchTodos(); }

  fetchTodos() { /* ... */ }
  handleTodoSubmit() { /* ... */ }
  handleDelete(id: number | undefined) { /* ... */ }
}
\end{lstlisting}

\begin{affianco}[0.44][c]
\codicepiccolo
\begin{lstlisting}[style=html, name=todo-page.component.html]
<ul>
  @for (item of todos; track item.id) {
    <app-todo-item [todoItem]="item"
                   (delete)="handleDelete($event)"/>
  } @empty {
    <li>Nessun elemento da mostrare.</li>
  }
</ul>
\end{lstlisting}
\risultato
\begin{immagine}
\begin{tikzpicture}
  \node[wtnode, font=\FiraCodeNL\scriptsize, minimum width=8cm] (p) at (0,0)
    {TodoPageComponent};
  \foreach \x [count=\i] in {-2.75,0,2.75}
    \node[wtnodeacc, font=\FiraCodeNL\scriptsize, inner sep=3pt] (f\i) at (\x,-2.1)
      {<app-todo-item>};
  \draw[wtarrow] (-3.05,-0.3) -- node[mcod, left]{[todoItem]} (-3.05,-1.85);
  \draw[wtarrowacc] (-2.45,-1.85) -- node[mcod, right]{(delete)} (-2.45,-0.3);
  \node[mnota, anchor=north] at (0,-2.45) {uno per ogni elemento di \code{todos}};
\end{tikzpicture}
\end{immagine}
\end{affianco}

\begin{lstlisting}[style=ts, name=todo-item.component.ts]
export class TodoItemComponent {
  @Input({ required: true }) todoItem: TodoItem;
  @Output() delete: EventEmitter<number | undefined> = new EventEmitter();

  editLink = "";
  restBackend = inject(RestBackendService);
  toastr      = inject(ToastrService);

  ngOnInit() { this.editLink = "/todos/" + this.todoItem?.id; }

  handleToggleDoneStatus() { /* ... */ }
  handleDelete() { /* ... */ }
}
\end{lstlisting}

\code{<app-todo-item>} è un componente a sé, responsabile di mostrare un singolo
elemento: riceve il dato tramite \code{@Input} e comunica al padre la richiesta
di cancellazione tramite \code{@Output}. È esattamente lo schema di
comunicazione descritto nel capitolo precedente, applicato a un caso reale.

\info{Il percorso completo}{
  Vale la pena fermarsi un momento a osservare che cosa abbiamo costruito: un
  backend Express con ORM che espone una API REST protetta da JWT (capitoli
  13--15), e un frontend Angular che la consuma con servizi iniettati, signal,
  guardie, interceptor e componenti (capitoli 19--20). Sono le stesse
  funzionalità della lista di cose da fare scritta in Bash e CGI nel capitolo 9:
  è cambiato tutto tranne il problema.}

\section{Riferimenti}

\begin{itemize}
  \item \textbf{Angular Docs}. \\
        \urlx{https://angular.dev/overview} \\
        Parti rilevanti: \textit{Introduction: Essentials}; \textit{In-depth
        Guides: Forms, Dependency Injection, Router, HTTP Client, Signals}.
\end{itemize}
